% GENERATED FILE. Edit canonical lessons, then run npm run generate:books.
% canonical-source-hash: fnv1a64:def9d135a0c47650
% canonical-lessons: ML-C82-chettaa, ML-C82-chechi-address, ML-C82-greeting-register, ML-R82-address-recall

\chapter{Opening A Conversation}
\label{ch:ml-address}
\hlchaptermodality{82}

\begin{chapteropening}
\textbf{By the end of this chapter:} \emph{I can call out to somebody I have never met, and I know which of my greetings to use with them.}

From cold: both calling words, the ending that makes one of them and why the other has nothing to change, and the rule for choosing among the five greetings already owned.
\end{chapteropening}

\section[cēṭṭā]{\ml{ചേട്ടാ} (cēṭṭā) — excuse me, brother}

\label{lesson:ML-C82-chettaa}



\begin{quote}
Say the word for an \textbf{older brother}. Then ask yourself something you have had no way to do yet: \textbf{how would you get a stranger's attention?}
\end{quote}

\subsection*{You'll want to know: \ml{ചേട്ടാ}}
\textbf{\ml{ചേട്ടാ}} — \emph{cēṭṭā} — the way you call out to a man you do not know.

You have owned \textbf{\ml{ചേട്ടൻ}}, \emph{older brother}, since the family words. Two things happen to it here, and they are worth separating.

\textbf{The first is who it is for.} In Kerala the sibling words are used well outside the family: a shopkeeper, a driver, somebody a little older than you in a queue. \textbf{You address a stranger as kin}, and doing so is ordinary politeness rather than a familiarity you have to earn. English has nothing that works this way — \emph{excuse me} is the nearest thing and it names no relationship at all.

\begin{grammarlens}[title={Grammar Lens: the calling form}]
\textbf{The second is the ending.} You do not call somebody with the same shape you use to talk \emph{about} them:

\noindent
\begin{tabularx}{\linewidth}{@{}>{\raggedright\arraybackslash}X>{\raggedright\arraybackslash}X@{}}
\toprule
\textbf{} & \textbf{} \\
\midrule
\textbf{\ml{ചേട്ടൻ}} \emph{cēṭṭan} & an older brother — the word for him \\
\textbf{\ml{ചേട്ടാ}} \emph{cēṭṭā} & \emph{brother!} — the word \textbf{to} him \\
\bottomrule
\end{tabularx}

\textbf{A word ending in -\ml{ൻ} swaps that ending for -\ml{ാ} when you call somebody with it.} That is the whole rule, and it reaches every such word you own: the one for a younger brother does the same thing, and so does the word for a teacher.

Malayalam keeps a separate shape for being spoken \emph{to}, where English changes nothing at all and relies on tone of voice.
\end{grammarlens}

\subsection*{Your turn}
Say these aloud:

\begin{itemize}
\raggedright
  \item \emph{cēṭṭan}, then \emph{cēṭṭā}
  \item which of the two you would use to get somebody's attention
  \item what happened to the ending
  \item the same change on the word for a younger brother
\end{itemize}

\begin{tcolorbox}[breakable,colback=teal!4,colframe=teal!35!black,title={Before you move on}]
What does \emph{cēṭṭā} do? (\textbf{Calls out to a man you do not know}.) Which ending does it replace? (\textbf{-\ml{ൻ}}.) And does the man have to be your brother? (\textbf{No — you call a stranger by the family word}.)
\end{tcolorbox}
\section[cēcci]{\ml{ചേച്ചി} (cēcci) — excuse me, sister}

\label{lesson:ML-C82-chechi-address}



\begin{quote}
Say the calling form for a man you do not know, and say which ending it replaced.
\end{quote}

\begin{grammarlens}[title={Grammar Lens: the ending that has nothing to change}]
\textbf{\ml{ചേച്ചി}} — \emph{cēcci} — the way you call out to a woman you do not know.

The move is the same one: \textbf{the word for an older sister, used for a stranger.} A shopkeeper, somebody at the next desk, a woman a little older than you. It is the partner of the word you learned a moment ago and it covers the same ground.

What is different is that \textbf{this one does not change shape.}

\noindent
\begin{tabularx}{\linewidth}{@{}>{\raggedright\arraybackslash}X>{\raggedright\arraybackslash}X>{\raggedright\arraybackslash}X@{}}
\toprule
\textbf{} & \textbf{for him / her} & \textbf{calling} \\
\midrule
older brother & \textbf{\ml{ചേട്ടൻ}} & \textbf{\ml{ചേട്ടാ}} \\
older sister & \textbf{\ml{ചേച്ചി}} & \textbf{\ml{ചേച്ചി}} \\
\bottomrule
\end{tabularx}

The \textbf{-\ml{ാ}} ending from the lesson before replaces \textbf{-\ml{ൻ}}, and \emph{cēcci} has no \textbf{-\ml{ൻ}} to replace. So it is called as it stands, usually with the last sound drawn out a little, which is a matter of voice rather than of spelling.

\textbf{A rule that applies to one word of a pair and not the other is normal}, and it is worth getting used to early. Nothing is broken here: this word carries no \textbf{-\ml{ൻ}} for the ending to take.
\end{grammarlens}

\subsection*{Your turn}
Say these aloud:

\begin{itemize}
\raggedright
  \item \emph{cēcci}
  \item \emph{cēṭṭā}, then \emph{cēcci}
  \item which of the two changed its ending, and why the other did not
\end{itemize}

\begin{tcolorbox}[breakable,colback=teal!4,colframe=teal!35!black,title={Before you move on}]
How do you call out to a woman you do not know? (\emph{\textbf{cēcci}}.) Does it take the calling ending? (\textbf{No — it has no -\ml{ൻ} to replace}.)
\end{tcolorbox}
\section[which greeting, and when]{which greeting, and when — five you own, one rule you do not}

\label{lesson:ML-C82-greeting-register}



\begin{quote}
Say the greeting you learned first, then the one for the morning.
\end{quote}

\begin{grammarlens}[title={Grammar Lens: the rule nobody told you}]
You own \textbf{five} greetings and no rule for choosing. A learner who picks wrong is not misunderstood — they sound like a radio announcer greeting a neighbour.

\noindent
\begin{tabularx}{\linewidth}{@{}>{\raggedright\arraybackslash}X>{\raggedright\arraybackslash}X@{}}
\toprule
\textbf{greeting} & \textbf{reach} \\
\midrule
\textbf{\ml{നമസ്കാരം}} \emph{namaskāraṁ} & \textbf{anyone, any hour} \\
\textbf{\ml{സുപ്രഭാതം}} \emph{suprabhātaṁ} & the morning, \textbf{formal or not} \\
\textbf{\ml{ശുഭ} \ml{മധ്യാഹ്നം}} \emph{śubha madhyāhnaṁ} & the afternoon, \textbf{formal} \\
\textbf{\ml{ശുഭ} \ml{സായാഹ്നം}} \emph{śubha sāyāhnaṁ} & the evening, \textbf{formal} \\
\textbf{\ml{ശുഭ} \ml{രാത്രി}} \emph{śubha rātri} & the night, \textbf{formal} \\
\bottomrule
\end{tabularx}

\textbf{\ml{നമസ്കാരം} is the one to reach for.} It is tied to no hour and no occasion, it is respectful without being stiff, and it is never wrong. If you remember one line from this page, that is the line.

\textbf{The \ml{ശുഭ} family is written and announced more than it is spoken.} These are the greetings of a notice, a broadcast, a speech, the opening of an email — you will read and hear them constantly, and say them rarely.

\textbf{\ml{സുപ്രഭാതം} is the exception inside its own family.} It is built on a different piece of Sanskrit than its neighbours and has ended up as the ordinary morning greeting rather than the ceremonial one. \textbf{Say it to anybody before noon.}

So the working version is two lines:

\begin{itemize}
\raggedright
  \item \textbf{Any time, anyone $\to$ \ml{നമസ്കാരം}.}
  \item \textbf{Before noon, if you like $\to$ \ml{സുപ്രഭാതം}.}
\end{itemize}

And the others are for reading.
\end{grammarlens}

\begin{grammarlens}[title={What you've built}]
Nothing new, and something you were missing.

\textbf{This lesson taught no word.} All five were already yours; what was missing was the thing a phrasebook never prints — \textbf{which of them a person actually says.}

Put it with the calling words and you can start a conversation from nothing: \textbf{\ml{നമസ്കാരം}, \ml{ചേട്ടാ}.}
\end{grammarlens}

\subsection*{Your turn}
Say these aloud:

\begin{itemize}
\raggedright
  \item the greeting that works at any hour with anyone
  \item what you would greet a neighbour with at nine in the morning
  \item which of the five you are most likely to READ rather than say
  \item \emph{namaskāraṁ, cēṭṭā}
\end{itemize}

\begin{tcolorbox}[breakable,colback=teal!4,colframe=teal!35!black,title={Before you move on}]
Which greeting is never wrong? (\emph{\textbf{namaskāraṁ}}.) Which family is mostly written and announced? (\textbf{The \ml{ശുഭ} one}.) And which member of that family is an everyday word anyway? (\emph{\textbf{suprabhātaṁ}}.)
\end{tcolorbox}
\section[sambhāṣaṇaṁ tuṭaṅṅal enna ōrmma]{(opening a conversation) — from cold}

\label{lesson:ML-R82-address-recall}



\begin{quote}
Close the lessons before this one. Call out to a man you do not know, then to a woman you do not know.
\end{quote}

\begin{grammarlens}[title={Grammar Lens: two things a stranger needs}]
\noindent
\begin{tabularx}{\linewidth}{@{}>{\raggedright\arraybackslash}X>{\raggedright\arraybackslash}X>{\raggedright\arraybackslash}X@{}}
\toprule
\textbf{} & \textbf{about them} & \textbf{\textbf{to} them} \\
\midrule
older brother & \textbf{\ml{ചേട്ടൻ}} & \textbf{\ml{ചേട്ടാ}} \\
older sister & \textbf{\ml{ചേച്ചി}} & \textbf{\ml{ചേച്ചി}} \\
\bottomrule
\end{tabularx}

\textbf{A word ending in -\ml{ൻ} swaps it for -\ml{ാ} when you call somebody}, and a word with no \textbf{-\ml{ൻ}} has nothing to swap, so it stands as it is.

Neither of these is a stranger-word that Malayalam keeps in reserve. \textbf{They are the family words}, used outward: you address somebody you have never met as an older brother or an older sister, and that is ordinary politeness rather than a familiarity you have to earn. English has no equivalent — \emph{excuse me} names no relationship at all.
\end{grammarlens}

\begin{grammarlens}[title={What you've built}]
An opening line, which is the one thing a learner needs first and usually gets last.

\textbf{\ml{നമസ്കാരം}, \ml{ചേട്ടാ}.}

Two pieces, and \textbf{one of them was free} — the greeting was learned in the very first chapter, and this chapter only said \textbf{when to use it}: any hour, anyone, never wrong. The four ceremonial greetings are for reading, and the morning one is the exception you can say to anybody.

\textbf{Knowing a word and knowing when it is used are two different pieces of knowledge.} A book that gives the first and not the second has given you half, and this chapter went back for the other half rather than adding more words.
\end{grammarlens}

\subsection*{Your turn}
Say these aloud:

\begin{itemize}
\raggedright
  \item \emph{cēṭṭā}, then \emph{cēcci}
  \item which of the two changed its ending, and what it changed from
  \item the greeting that is never wrong
  \item \emph{namaskāraṁ, cēṭṭā}
  \item \emph{namaskāraṁ, cēcci}
\end{itemize}

\begin{tcolorbox}[breakable,colback=teal!4,colframe=teal!35!black,title={Before you move on}]
Which ending does the calling form replace? (\textbf{-\ml{ൻ}, with -\ml{ാ}}.) Whose words are you using when you address a stranger? (\textbf{The family's}.) And which greeting can you always fall back on? (\emph{\textbf{namaskāraṁ}}.)
\end{tcolorbox}
